\chapter{2025年天津卷}

\section{单选题}

\begin{problem}
已知集合 \(\mathcal U=\{1,2,3,4,5\}\),\(A=\{1,3\}\),\(B=\{2,3,5\}\),则 \(\mathcal U\setminus(A\cup B)=\)(\quad)

\choices
    {\(\{1,2,3,4\}\)}
    {\(\{2,3,4\}\)}
    {\(\{2,4\}\)}
    {\(\{4\}\)}
\end{problem}

\begin{problem}
设 \(x\in\R\),则``\(x=0\)''是``\(\sin 2x=0\)''的(\quad)

\choices
    {充分不必要条件}
    {必要不充分条件}
    {充分必要条件}
    {既不充分也不必要条件}
\end{problem}

\begin{problem}
已知函数 \(y=f(x)\) 的图象如下,则 \(f(x)\) 的解析式可能为(\quad)

\bitmapfigure[max width=0.42\linewidth]{img/2025/tianjin/q3_fig1.png}

\choices
    {\(f(x)=\frac{x}{1-|x|}\)}
    {\(f(x)=\frac{x}{|x|-1}\)}
    {\(f(x)=\frac{|x|}{1-x^2}\)}
    {\(f(x)=\frac{|x|}{x^2-1}\)}
\end{problem}

\begin{problem}
若 \(m\) 为直线,\(\alpha\),\(\beta\) 为两个平面,则下列结论中正确的是(\quad)

\choices
    {若 \(m\parallel \alpha\),\(n\subset \alpha\),则 \(m\parallel n\)}
    {若 \(m\perp\alpha\),\(m\perp\beta\),则 \(\alpha\perp\beta\)}
    {若 \(m\parallel\alpha\),\(m\perp\beta\),则 \(\alpha\perp\beta\)}
    {若 \(m\subset\alpha\),\(\alpha\perp\beta\),则 \(m\perp\beta\)}
\end{problem}

\begin{problem}
下列说法中错误的是(\quad)

\choices
    {若 \(X\sim N(\mu,\sigma^2)\),则 \(P(X\le\mu-\sigma)=P(X\ge\mu+\sigma)\)}
    {若 \(X\sim N(1,2^2)\),\(Y\sim N(2,2^2)\),则 \(P(X<1)<P(Y<2)\)}
    {\(|r|\) 越接近 1,相关性越强}
    {\(|r|\) 越接近 0,相关性越弱}
\end{problem}

\begin{problem}
\(S_n=-n^2+8n\),则数列 \(\{|a_n|\}\) 的前 12 项和为(\quad)

\choices
    {112}
    {48}
    {80}
    {64}
\end{problem}

\begin{problem}
函数 \(f(x)=0.3^x-\sqrt{x}\) 的零点所在区间是(\quad)

\choices
    {\((0,0.3)\)}
    {\((0.3,0.5)\)}
    {\((0.5,1)\)}
    {\((1,2)\)}
\end{problem}

\begin{problem}
\(f(x)=\sin(\omega x+\varphi)\)(\(\omega>0\),\(-\pi<\varphi<\pi\)),在 \(\left[-\frac{5\pi}{12},\frac{\pi}{12}\right]\) 上单调递增,且 \(x=\frac{\pi}{12}\) 是它的一条对称轴,\((\frac{\pi}{3},0)\) 是它的一个对称中心.

当 \(x\in[0,\frac{\pi}{2}]\) 时,\(f(x)\) 的最小值为(\quad)

\choices
    {\(-\frac{\sqrt3}{2}\)}
    {\(-\frac12\)}
    {1}
    {0}
\end{problem}

\begin{problem}
双曲线 \(\frac{x^2}{a^2}-\frac{y^2}{b^2}=1\)(\(a>0\),\(b>0\))的左,右焦点分别为 \(F_1\),\(F_2\),以右焦点 \(F_2\) 为焦点的抛物线 \(y^2=2px\)(\(p>0\))与双曲线在第一象限交于 \(P\),若 \(|PF_1|+|PF_2|=3|F_1F_2|\),则双曲线离心率 \(e=\)(\quad)

\choices
    {2}
    {5}
    {\(\frac{\sqrt2+1}{2}\)}
    {\(\frac{\sqrt5+1}{2}\)}
\end{problem}

\section{填空题}

\begin{problem}
已知 \(\i\) 是虚数单位,则 \(\left|\frac{3+\i}{\i}\right|=\fillinblank{}\).
\end{problem}

\begin{problem}
在 \((x-1)^6\) 的展开式中,\(x^3\) 项系数为 \(\fillinblank{}\).
\end{problem}

\begin{problem}
\(l_1:x-y+6=0\) 与 \(x\) 轴,\(y\) 轴交点分别为 \(A\),\(B\),又与 \((x+1)^2+(y-3)^2=r^2\) 交于 \(C\),\(D\),且 \(|AB|=3|CD|\),则 \(r=\fillinblank{}\).
\end{problem}

\begin{problem}
小桐一周跑 \(2\) 次,一次 \(5\) 圈或 \(6\) 圈. 第一次跑 \(5\) 圈或 \(6\) 圈概率各为 \(0.5\);若第一次跑 \(5\) 圈,则第二次分别跑 \(5\) 圈,\(6\) 圈的概率为 \(0.4\),\(0.6\);若第一次跑 \(6\) 圈,则第二次分别跑 \(5\) 圈,\(6\) 圈的概率为 \(0.6\),\(0.4\).
则一周跑 \(11\) 圈的概率为 \(\fillinblank{}\);若一周至少跑 \(11\) 圈为达标,则连续跑 \(4\) 周,记合格周数为 \(X\),则 \(E(X)=\fillinblank{}\).
\end{problem}

\begin{problem}
\(\triangle ABC\) 中,\(D\) 为 \(AB\) 中点,\(\overrightarrow{CE}=\frac13\overrightarrow{CD}\),\(\overrightarrow{AC}= \bs b\),\(\overrightarrow{AB}=\bs a\).
则 \(\overrightarrow{AE}=\fillinblank{}\)(用 \(\bs a\),\(\bs b\) 表示);若 \(|\overrightarrow{AE}|=5\),\(\overrightarrow{AE}\perp\overrightarrow{CB}\),则 \(\overrightarrow{AE}\cdot\overrightarrow{CD}=\fillinblank{}\).

\FigureLayoutDeclare{img/2025/tianjin/q14_fig1.png}{5.5cm}{0bp 0bp 0bp 0bp}
\bitmapfigure[max width=0.42\linewidth]{img/2025/tianjin/q14_fig1.png}
\end{problem}

\begin{problem}
若 \(a\),\(b\in\R\),对任意 \(x\in[-2,2]\),都有 \((2a+b)x^2+bx-a-1\le0\),则 \(2a+b\) 的最小值为 \(\fillinblank{}\).
\end{problem}

\section{解答题}

\begin{problem}
在 \(\triangle ABC\) 中,角 \(A\),\(B\),\(C\) 对边分别为 \(a\),\(b\),\(c\). 已知 \(a\sin B=\sqrt3\,b\cos A\),\(c-2b=1\),\(a=\sqrt7\).
\begin{enumerate}
    \item 求 \(A\);
    \item 求 \(c\);
    \item 求 \(\sin(A+2B)\).
\end{enumerate}
\end{problem}

\begin{problem}
正方体 \(ABCD-A_1B_1C_1D_1\) 的棱长为 4,\(E\),\(F\) 分别为 \(A_1D_1\),\(C_1B_1\) 中点,且 \(CG=3GC_1\).

\begin{enumerate}
\item 求证:\(GF\perp\) 平面 \(FBE\);
\item 求平面 \(FBE\) 与平面 \(EBG\) 的夹角余弦;
\item 求四面体 \(D-FBE\) 的体积.
\end{enumerate}

\bitmapfigure[max width=0.55\linewidth]{img/2025/tianjin/q17_fig1.png}
\end{problem}

\begin{problem}
已知椭圆 \(\frac{x^2}{a^2}+\frac{y^2}{b^2}=1\)(\(a>b>0\))的左焦点为 \(F\),右顶点为 \(A\),\(P\) 为 \(x=a\) 上一点,且直线 \(PF\) 的斜率为 \(\frac13\),\(\triangle PFA\) 的面积为 \(\frac32\),离心率为 \(\frac12\).
\begin{enumerate}
\item 求椭圆方程;
\item 过 \(P\) 的直线与椭圆有唯一交点 \(B\)(异于 \(A\)),证明 \(PF\) 平分 \(\angle AFB\).
\end{enumerate}

\bitmapfigure[max width=0.42\linewidth]{img/2025/tianjin/q18_fig1.png}
\end{problem}

\begin{problem}
已知数列 \(\{a_n\}\) 为等差数列,\(\{b_n\}\) 为等比数列,\(a_1=b_1=2\),\(a_2=b_2+1\),\(a_3=b_3\).
\begin{enumerate}
\item 求 \(\{a_n\}\),\(\{b_n\}\) 的通项;
\item \(\forall n\in\N^*\),\(I=\{0,1\}\),定义 \(T_n=\{p_1a_1b_1+p_2a_2b_2+\cdots+p_na_nb_n\mid p_i\in I,i=1,\dots,n\}\).
\begin{enumerate}
\item 求证:对任意 \(t\in T_n\),有 \(t<a_{n+1}b_{n+1}\);
\item 求 \(T_n\) 的所有元素之和.
\end{enumerate}
\end{enumerate}
\end{problem}

\begin{problem}
已知 \(f(x)=ax-(\ln x)^2\).
\begin{enumerate}
\item 当 \(a=1\) 时,求 \(f(x)\) 在点 \((1,f(1))\) 处的切线方程;
\item 当 \(f(x)\) 有三个零点 \(x_1<x_2<x_3\) 时,
\begin{enumerate}
    \item 求 \(a\) 的取值范围;
    \item 证明 \((\ln x_2-\ln x_1)\cdot \ln x_3<\frac{4\e}{\e-1}\).
\end{enumerate}
\end{enumerate}
\end{problem}

